% ============================================================================
%  Part I: Why a New Identifier
% ============================================================================
\section{Why a New Identifier}
\label{sec:why}

Every identifier format reflects the constraints and assumptions of the moment it was designed. The decisions baked into UUIDv4 \citep{rfc4122} in 2005 --- 128 bits total, 122 bits of entropy, no time component, no ordering --- made sense when distributed databases were rare, quantum computers were a footnote, and ``unique'' meant ``unique in this rack.'' Twenty years later, the constraints have shifted enough that even the carefully-chosen successors to UUIDv4 (UUIDv7 \citep{rfc9562}, ULID \citep{ulid}, KSUID \citep{ksuid}) leave gaps. PQID is an attempt to fill the remaining gaps in a single, coherent format --- not by adding features, but by making slightly different choices on three numbers: the size of the random field, the granularity of the timestamp, and the semantics of the version byte.

\subsection{Quantum search has changed the entropy budget}

Grover's algorithm \citep{grover1996} gives a quantum computer a quadratic speedup over classical search: finding a single matching value among $2^n$ candidates takes approximately $2^{n/2}$ quantum operations. For an identifier whose security model rests on the difficulty of guessing a specific value (a \emph{targeted second preimage}), that means an $n$-bit random field gives only $n/2$ bits of effective post-quantum security. UUIDv4's 122 bits give 61. UUIDv7's 74 give 37. ULID's 80 give 40. KSUID's 128 give 64. PQID's 256 give 128 --- the level NIST has designated as adequate for symmetric primitives through 2030 and beyond \citep{nist80057}, and the level ML-DSA-44 \citep{fips204} and ML-KEM-512 \citep{fips203} are calibrated to. Choosing less for a new identifier means accepting that the identifier becomes the weakest cryptographic component of any system that combines it with AES-128 \citep{fips197} or stronger. Part \ref{sec:security} works through the math (Grover, BHT \citep{bht1998}, the classical and quantum birthday bounds) in detail.

\subsection{Time-sortability is a feature, not a side effect}

Random identifiers like UUIDv4 are catastrophic for B-tree indexes. Each new insert lands at a random point in the tree, causing scattered page writes, poor cache behavior, and eventually fragmentation that degrades read performance. The fix is well-understood: prefix the identifier with a timestamp so consecutive inserts land near each other. ULID, KSUID, and UUIDv7 all do this. PQID does too.

But there's a subtlety. ULID and KSUID expose lexicographic byte ordering as a first-class property: sorting their bytes (or, in ULID's case, their Crockford-encoded strings) yields time-ordering. UUIDv7's RFC permits implementations that randomize within a millisecond in ways that break this property; the spec only guarantees ordering \emph{between} milliseconds, not \emph{within} them. PQID guarantees both. Within a single millisecond, a per-generator counter makes consecutive IDs strictly increasing; across milliseconds, the timestamp dominates. The result is 100\% time-sortedness across consecutive IDs from a single generator (binary order, and ASCII order in the Crockford encoding), measured against 10{,}000 consecutive IDs in the benchmark suite and against 100{,}000 unique IDs in the uniqueness test.

\subsection{Distributed systems need first-class causal ordering}

The Standard variant uses a wall-clock timestamp. That works fine when a single machine is producing IDs --- the local clock is the source of truth and consumers that care about ordering get what they expect. It works less well when many machines are producing IDs and clocks disagree by tens of milliseconds.

The HLC variant solves this. A Hybrid Logical Clock \citep{kulkarni2014hlc} combines physical time with a logical counter. An HLC reading is always at least the maximum of the local wall clock and any HLC value the node has observed from a peer; the counter increments whenever logical time fails to advance. The result is a clock that stays close to physical time, always moves forward, and preserves causal ordering across nodes that exchange HLC values in their messages.

PQID's HLC variant exposes that machinery as a first-class identifier shape. The wire format is identical to the Standard variant; the only difference is the variant byte (\code{0x02} instead of \code{0x01}) and the semantic interpretation of the timestamp (logical, not physical, with a 60-second drift bound). A consumer that doesn't care about variant semantics can sort PQIDs as opaque bytes and get correct ordering across both variants in the same stream.

\subsection{The cost: bytes}

PQID's distinguishing properties are the 256-bit random field --- the only format with comfortable post-quantum margins --- and the dual-variant design that makes the same wire format usable for both centralized and distributed generation. The cost is byte-count: at 41 bytes binary or 55 characters Base64URL, PQID is roughly $2.5\times$ the size of UUIDv4. Part \ref{sec:compare} works through the format-by-format tradeoffs; Part \ref{sec:when} maps them onto specific use cases.

